Zur Auswahl der Lösungmethoden für die verwendete Mechanik wurde wie folgt vorgegangen:
\begin{enumerate}
    \item Extraktion der mechanischen Anforderungen aus dem Lastenheft
    \item Gruppieren der notwendigen Lösungsdimensionen
    \item Explorative Lösungsfindung mit Hilfe einer MindMap
    \item Gegenüberstellung der verschiedenen Lösungswege
    \item Begründung der verfolgten Lösung
\end{enumerate}

Da die unterschiedlichen Lösungsoptionen und Gegenüberstellungen ausführlich im Dokument "Alternativenbewertung"
beschrieben wurden und diesem im DMS der WBH zur Verfügung steht, werden an dieser Stelle lediglich die Gruppierun
und die gewählten Lösungswege mit ihren Begründungen vorgestellt. 
Die Gruppierung der Anforderungen für die Mechnik sind wie folgt:
\begin{itemize}
    \item Chassis
    \item[] \begin{itemize}
        \item Material
        \item Form
        \item Modularität
    \end{itemize}
    \item Untersetzer ablegen
    \item Untersetzer aufnehmen
\end{itemize} 

\textbf{Chassis - Materiel}
Als Material für das Chassis des Coasterbots fiel die Wahl auf Kunststoff (Polylactid, PLA). Dieses Material
wurde ausgewählt, da es als günstig, besonders leicht und einfach zu verarbeiten mit Hilfe eines 3D-Druckers
gilt. 